\section{拘留}
\begin{enumerate}
  \item 区分行政拘留、司法拘留。公安执行。
  \item \zd{临时性突发性}；日做单位的；逮捕是相对更加稳定的一种羁押状态。
  \item 主体。
        \begin{enumerate}
          \item 检察院。犯罪后企图\bs{自杀、逃跑、在逃的；毁灭、伪造、串供的}。
          \item 公安。\bs{一言不合就带走}。这个人可疑，犯罪后逃避刑事责任。
                \begin{enumerate}
                  \item \bs{正在、指认、发现、证}。正在预备犯罪；指认犯罪；发现犯罪证据。
                  \item \bs{毁、逃、杀、身、流、结、多}。毁灭、伪造证据；企图逃跑、自杀；身份不明；流窜、结伙、多次作案重大嫌疑。
                \end{enumerate}
          \item \bs{法院拘留一定是司法拘留}。不是刑事强制措施。
        \end{enumerate}
  \item 程序。
        \begin{enumerate}
          \item 紧急情况经查证\zd{可以没有拘留证}，24 小时内\zd{补办手续}。
          \item \bs{24 小时内 3 件事。送看、通知、讯问}。送看守所、通知家属（危恐除外）、讯问。\al{黑社会性质犯罪 \cuo}。
                \al{送看守所后 24 小时内通知家属 \cuo}。
                \al{未通知家属、未告知家属理由 \cuo}。
          \item 异地拘留。应当通知当地公安协助\zd{配合}。
        \end{enumerate}
  \item 期限。
        \begin{enumerate}
          \item 公安。X 日内报批 + Y 日检决定。
                \begin{enumerate}
                  \item 一般 3/复杂 7/流结多 30 \yjt 流窜、结伙（2 人以上）、多次（3 次以上）作案。
                  \item \zd{3/7/30 + 7 = Max(10, 14, 37) = 37}。\zd{实践中 37 天没捞出来，喜提逮捕}。
                \end{enumerate}
          \item 检察院。\zd{7 + (7 + 3) = Max(17)}。
          \item 监委。\zd{10/14 = Max(14)}。
                \begin{enumerate}
                  \item 叫\zd{留置}（\bdy 强制措施） \yjt 移送审查起诉之日\zd{自动解除} \yjt 检\zd{先行拘留}。
                  \item 市以下的，应当报上一级监委批准。
                  \item \al{留置期间，检可以对监管活动进行监督 \cuo{监委强势，不得插手}}。
                  \item 管制 1:2 折抵，拘役、有期 1:1 折抵。
                \end{enumerate}
        \end{enumerate}
  \item 后续。逮捕、变更、释放。
\end{enumerate}
